The assertion is trivial when \(\sigma=0\), so first suppose \(\sigma>0\).  We work inside a subclass having independent \(N(0,\sigma^2)\) innovations across rounds and outcome coordinates, zero direct effects unless explicitly specified, and the additive response surfaces below.  This is a subclass of \(\mathcal M_s\): every constructed row has at most one spillover, every nonzero spillover has magnitude exactly \(\delta\), and Gaussian innovations satisfy the displayed conditional sub-Gaussian condition.

Define the rare-side indicator
\[
 a_j(z)=
 \begin{cases}
 z_j,&p\le1/2,\\
 1-z_j,&p>1/2.
 \end{cases}
\]
For every legal action, \(\sum_{j=1}^na_j(z)=nq\).  Let \(P_0\) have all response means zero and hence empty graph.  For each \(j\), choose a recipient \(i(j)\ne j\), and define \(P_j\) so that its only nonzero mean coordinate is
\[
 \mu_{i(j)}(z)=\delta a_j(z).
\]
When \(p\le1/2\), this is obtained with \(\Gamma_{i(j),j}=\delta\) and zero baseline.  When \(p>1/2\), it is obtained with \(\alpha_{i(j)}=\delta\) and \(\Gamma_{i(j),j}=-\delta\).  Thus \(P_j\) has exactly one directed edge.

The chain rule for likelihood ratios under an adaptive design has no assignment contribution, because the same predictable kernel \(\pi_t(\cdot\mid\mathcal F_{t-1})\) is used under both environments.  Conditional on a history and action, the Gaussian outcome divergence is the squared mean difference divided by \(2\sigma^2\).  Therefore
\[
 \operatorname{KL}(\Pr_{P_0,d},\Pr_{P_j,d})
 =\frac{\delta^2}{2\sigma^2}
   \mathbb E_{P_0,d}\sum_{t=1}^m a_j(Z_{t\cdot}).                      \tag{25}
\]
Summing the right side over \(j\) shows that some \(j_*\) satisfies
\[
 \operatorname{KL}(\Pr_{P_0,d},\Pr_{P_{j_*},d})
 \le\frac{m q\delta^2}{2\sigma^2}.                                   \tag{26}
\]
Let \(E_0=\{g(\mathscr H_m)=A(P_0)\}\).  Uniform graph error at most \(\eta\) gives
\[
 \Pr_{P_0,d}(E_0)\ge1-\eta,
 \qquad
 \Pr_{P_{j_*},d}(E_0)\le\eta.
\]
Applying the log-sum inequality to the two cells \(E_0,E_0^c\) yields the binary data-processing bound
\[
 \operatorname{KL}(\Pr_{P_0,d},\Pr_{P_{j_*},d})
 \ge (1-2\eta)\log\!\left(\frac{1-\eta}{\eta}\right)
 \ge c_1\log(1/\eta),                                                 \tag{27}
\]
where the last inequality uses \(0<\eta\le1/4\).  Equations (26)--(27) prove
\[
 m\ge c_2\frac{\sigma^2\log(1/\eta)}{q\delta^2}.                      \tag{28}
\]

For the \(\log(n-1)\) term, let \(a=\lfloor n/2\rfloor\), and choose disjoint sets \(I,J\subseteq\mathcal U\) with \(|I|=|J|=a\).  For each bijection \(h:I\to J\), define \(P_h\) by
\[
 \mu_i(z)=\delta a_{h(i)}(z)\quad(i\in I),
 \qquad
 \mu_i(z)=0\quad(i\notin I).
\]
Its graph consists of the \(a\) edges \((i,h(i))\), and every row has degree at most one.  For every legal action,
\[
 \lVert\mu^{(h)}(z)-\mu^{(0)}(z)\rVert_2^2
 =\delta^2\sum_{j\in J}a_j(z)
 \le\delta^2nq.
\]
The same adaptive chain-rule calculation, now in the orientation \(P_h\) versus \(P_0\), gives
\[
 \operatorname{KL}(\Pr_{P_h,d},\Pr_{P_0,d})
 \le\frac{mnq\delta^2}{2\sigma^2}.                                   \tag{29}
\]
Let \(H\) be uniform on the \(M=a!\) bijections and let \(\overline P=M^{-1}\sum_h\Pr_{P_h,d}\).  Direct expansion of the logarithms gives
\[
 \frac1M\sum_h\operatorname{KL}(\Pr_{P_h,d},\Pr_{P_0,d})
 =I(H;\mathscr H_m)+\operatorname{KL}(\overline P,\Pr_{P_0,d}),
\]
so (29) upper-bounds the mutual information.  On the other hand, the graph decoder determines \(h\) whenever it is correct.  If \(E\) is its error indicator, the chain rule for entropy gives
\[
 H(H\mid\mathscr H_m)
 \le H(E)+\Pr(E=1)\log(M-1)
 \le\log2+\eta\log M.
\]
Consequently,
\[
 I(H;\mathscr H_m)
 \ge(1-\eta)\log M-\log2
 \ge c_3\log M.                                                       \tag{30}
\]
Here \(a\ge4\), \(M\ge24\), and \(\eta\le1/4\).  Also, the last \(\lfloor a/2\rfloor\) factors in \(a!\) give
\[
 \log(a!)\ge\frac a2\log\!\left(\frac a2\right)
 \ge c_4n\log(n-1)\qquad(n\ge8).                                    \tag{31}
\]
Combining (29)--(31) yields
\[
 m\ge c_5\frac{\sigma^2\log(n-1)}{q\delta^2}.                        \tag{32}
\]
The maximum of (28) and (32) is at least half their sum, proving the first displayed lower bound after changing the universal constant.

It remains to prove the noiseless combinatorial assertion.  Interpret uniform identification here as probability-one correctness for every environment in the degree-one noiseless subclass.  Couple the design's internal randomization across the empty-graph environment and the finite family consisting of all one-edge alternatives used below.  The intersection of their probability-one correctness events has probability one.  Fix a seed in that intersection.  Under the all-zero environment, its all-zero history makes the adaptive rule produce a deterministic legal matrix \(Z\).

Every column of this matrix must be nonconstant.  If column \(j\) is identically zero, a coefficient \(\Gamma_{ij}=\delta\), with \(i\ne j\), is invisible along the all-zero path.  If it is identically one, the same edge with baseline \(\alpha_i=-\delta\) is invisible.  Distinct columns must also be distinct: if columns \(j\) and \(k\) agree, take recipient \(k\), set \(\Gamma_{k j}=\delta\) and \(\tau_k=-\delta\); then
\[
 \mu_k(Z_{t\cdot})=\delta Z_{tj}-\delta Z_{tk}=0
 \quad\text{for every }t.
\]
Each such alternative has one spillover edge, yet it generates exactly the all-zero history and hence the same terminal output as the empty graph, a contradiction.  Therefore the \(n\) binary column words, together with the all-zero and all-one words, are pairwise distinct.  Thus
\[
 2^m\ge n+2,
 \qquad
 m\ge\lceil\log_2(n+2)\rceil.                                       \tag{33}
\]
Nonconstancy also means that every column contains at least one one and at least one zero.  Counting entries by rows gives
\[
 mB\ge n,
 \qquad
 m(n-B)\ge n.
\]
Since \(q=\min\{B/n,1-B/n\}\), these inequalities imply
\[
 m\ge\max\left\{\left\lceil\frac nB\right\rceil,
                    \left\lceil\frac n{n-B}\right\rceil\right\}
 =\lceil1/q\rceil.                                                     \tag{34}
\]
Equations (33)--(34) prove the noiseless degree-one statement.